% chapters/glossary.tex — Glossary (in alphabetical order)

\chapter*{Glossary}
\addcontentsline{toc}{chapter}{Glossary}
\markboth{Glossary}{Glossary}

This glossary organizes the main terms used in this book in alphabetical order. Each item is organized with a short definition and the corresponding chapter number for easy reference by readers. When you first read the book, skim it lightly, and if you come across a term you don't know while reading the text, come back and refer to it.

\begin{description}[leftmargin=2.5em, style=nextline, font=\bfseries\sffamily\color{inkblue}]

\item [Activation (activation value)]
The value output after passing through each layer of the neural network. This refers to the output tensor of the layer and becomes the input to the next layer.\textit{reference: Chapter 4}

\item [AdamW]
An optimization algorithm that combines the Adam optimizer with weight decay. The de facto standard for modern LLM learning.\textit{reference: Chapter 7}

\item [Attention]
A mechanism by which each position in a sequence learns how much attention to pay to every other position. The core of Transformers.\textit{reference: Chapter 5}

\item [Backpropagation]
An algorithm that calculates the gradient by differentiating the loss of a neural network by each parameter. Based on the chain rule.\textit{reference: Chapter 7}

\item [BPE (Byte Pair Encoding)]
Tokenizer algorithm starts byte by byte and merges the most frequently occurring pairs. Used by GPT-2, GPT-3, etc.\textit{reference: Chapter 3}

\item [Causal Mask]
A lower triangle mask that blocks future tokens from being seen in an autoregressive model.\textit{reference: Chapter 5}

\item [Context Length]
Maximum number of tokens the model can process at one time. GPT-2 is 1024, GPT-4 is 128K.\textit{reference: Chapter 7}

\item [Cross-Entropy Loss]
A loss function that measures the difference between the predicted probability distribution and the correct answer distribution in a classification problem. Standard loss in language modeling.\textit{reference: Chapter 7}

\item [Decoder]
The part of the transformer that generates sequences. GPT is a structure that uses only a decoder.\textit{reference: Chapter 6}

\item [Dropout]
A regularization technique that prevents overfitting by randomly setting some neurons to 0 during training.\textit{reference: Chapter 4}

\item [Embedding]
The process of converting a discrete token ID into a continuous dense vector. Or the vector itself.\textit{reference: Chapter 4}

\item [Feed-Forward Network (FFN)]
Fully connected neural network for each position applied after attention in the transformer block. Usually 2 linear layers + activation function.\textit{reference: Chapter 6}

\item [GELU (Gaussian Error Linear Unit)]
A smooth version of ReLU. Smooth transition around zero using a Gaussian function. GPT model is used.\textit{reference: Chapter 6}

\item [Gradient]
Differential values ​​for each parameter of the loss function. Indicates the direction of optimization.\textit{reference: Chapter 7}

\item [Greedy Decoding]
A generation strategy that selects only the tokens with the highest probability at each step. Fast but monotonous output.\textit{reference: Chapter 8}

\item [Head]
A unit that independently performs attention in multihead attention.\textit{reference: Chapter 5}

\item [Layer Normalization]
A technique for normalizing the mean and variance along the last dimension of a tensor. Essential for stabilizing transformer learning.\textit{reference: Chapter 4, Chapter 6}

\item [Learning Rate]
Hyperparameter that determines the step size of the gradient. If it is too large, it will diverge; if it is too small, it will converge slowly.\textit{reference: Chapter 7}

\item [Logits]
Raw output values ​​before applying softmax. Usually the output of the last linear layer of the model.\textit{reference: Chapter 7}

\item [Loss Function]
A function that measures the difference between the model's predictions and the correct answer as a scalar. Optimization target for learning.\textit{reference: Chapter 7}

\item [Multi-Head Attention]
Learn various patterns simultaneously by running multiple attention heads in parallel.\textit{reference: Chapter 5}

\item [Parameter]
Weights that the model learns. W, b, etc. for linear layers.\textit{reference: Chapter 7}

\item [Perplexity]
Evaluation metrics for language models.$PPL = \exp(\text{loss})$. The lower the model, the better.\textit{reference: Chapter 7}

\item [Positional Encoding]
How to inject ordering information into a transformer. sin/cos or learnable embeddings.\textit{reference: Chapter 4}

\item [Pre-LN / Post-LN]
Location of layer normalization. Pre-LN is applied before attention/FFN, and Post-LN is applied after. Pre-LN has excellent learning stability.\textit{reference: Chapter 6}

\item [Residual Connection]
A connection that adds an input to an output.$y = x + f(x)$. Essential for deep network training.\textit{reference: Chapter 6}

\item [Scaled Dot-Product Attention]
Attention in the form$Q \cdot K^T / \sqrt{d\_k}$. The most basic operation of transformers.\textit{reference: Chapter 5}

\item [Self-Attention]
Attention between tokens in the same sequence. The core of Transformers.\textit{reference: Chapter 5}

\item [Softmax]
A function that converts a vector into a probability distribution.$\text{softmax}(x\_i) = e^{x\_i} / \sum\_j e^{x\_j}$.\textit{reference: Chapter 5}

\item [Temperature]
A parameter that controls the sharpness of the distribution when sampling. Lower is deterministic, higher is random.\textit{reference: Chapter 8}

\item [Token]
The smallest unit by which text is divided. Can be a word, subword, or byte.\textit{reference: Chapter 3}

\item [Tokenizer]
A module that converts text into token ID sequences and vice versa.\textit{reference: Chapter 3}

\item [Top-K Sampling]
A sampling strategy that considers only the K items with the highest probability among the following token candidates.\textit{reference: Chapter 8}

\item [Top-P (Nucleus) Sampling]
A sampling strategy that considers only tokens until the cumulative probability reaches P. Adaptive to distribution type.\textit{reference: Chapter 8}

\item [Transformer]
Vaswani et al. Attention-based neural network structure proposed by (2017). The foundation of the modern LLM.\textit{reference: Chapter 6}

\item [Unigram Tokenizer]
A tokenizer that assigns a probability to each token and learns with EM. Used by SentencePiece.\textit{reference: Chapter 3}

\item [Viterbi Algorithm]
Algorithm for finding optimal partitioning using dynamic programming. Used for decoding of Unigram tokenizer.\textit{reference: Chapter 3}

\item [Vocabulary]
The set of all tokens recognized by the tokenizer. Vocabulary size is usually 10,000 to 100,000.\textit{reference: Chapter 3}

\item [Weight Tying]
A technique to reduce the number of parameters by sharing the weights of the input embedding and output layer.\textit{reference: Chapter 7}

\end{description}
